% =============================================================
% Bölüm 3: Kapı Seviyesinde Sadeleştirme (K-map)
% =============================================================
\chapter{Kapı Seviyesinde Sadeleştirme (K-map)}
\label{chap:kmap}

\begin{tcolorbox}[
    colback=blokyesil!15,
    colframe=sinyalyesil!60!black,
    title={\textbf{Öğrenme Hedefleri}},
    fonttitle=\bfseries,
    rounded corners,
    boxrule=0.8pt
]
Bu bölümü tamamladığınızda aşağıdakileri yapabileceksiniz:
\begin{itemize}[nosep, leftmargin=*]
    \item 2, 3 ve 4 değişkenli Karnaugh haritalarını (K-map) çizmek ve bir Bool
    fonksiyonunu haritada göstermek.
    \item Prime implicant ve essential prime implicant kavramlarını kullanarak bir
    fonksiyonun minimum toplamlar-çarpımı (SOP) ve çarpımlar-toplamı (POS) biçimlerini
    haritadan doğrudan okumak.
    \item Bir doğruluk tablosundan K-map oluşturmak.
    \item Önemsenmeyen (don't-care) durumları sadeleştirme sürecinde kullanmak.
    \item Bir fonksiyonu yalnızca NAND ya da yalnızca NOR kapılarıyla iki-seviyeli
    gerçeklemek.
    \item XOR/XNOR fonksiyonlarının K-map'teki karakteristik (çapraz) desenini tanımak.
    \item Quine-McCluskey algoritmasıyla, elle harita çizmenin pratik olmadığı çok
    değişkenli fonksiyonları sistematik/tablo yöntemiyle sadeleştirmek.
\end{itemize}
\end{tcolorbox}

\section{Giriş}
\label{sec:kmap-giris}

Önceki bölümde Bool cebirinin aksiyomlarını ve teoremlerini kullanarak bir ifadeyi elle
sadeleştirmeyi öğrendik: De Morgan teoremiyle tümleyen almak, konsensüs teoremiyle fazladan
terimi atmak, dağılma/yutma gibi teoremlerle terimleri birleştirmek. Bu yöntem doğru sonucu
verir, ama önemli bir eksiği vardır: hangi teoremin hangi adımda uygulanacağını söyleyen
sabit bir kural yoktur. İki mühendis aynı fonksiyonu farklı sıralarda sadeleştirmeye
çalışıp farklı ara adımlardan geçebilir; hatta bazen hangi adımın atılacağını görmek
neredeyse bir sezgi/deneme-yanılma işine döner. Az sayıda değişkenli küçük bir ifadede bu
sorun olmaz, ama gerçek devrelerde onlarca terimli ifadelerle karşılaşıldığında, adım adım
işleyen açık bir prosedürün eksikliği ciddi bir engele dönüşür.

Bu bölümün konusu olan \textbf{Karnaugh haritası}\index{Karnaugh haritası} (kısaca
\textbf{K-map}), tam olarak bu boşluğu dolduran, tamamen görsel/sistematik bir sadeleştirme
yöntemidir.

\begin{figure}[H]
\centering
\includegraphics[height=4cm,keepaspectratio]{maurice_karnaugh.jpg}
\caption{Maurice Karnaugh (1924--2022).}
\label{fig:maurice-karnaugh}
\end{figure}

Yöntem, adını 1924'te New York'ta doğan Amerikalı fizikçi-mühendis \textbf{Maurice
Karnaugh}'dan alır. Yale Üniversitesi'nde fizik doktorasını tamamladıktan sonra 1952'de
Bell Laboratuvarları'na katılan Karnaugh, telefon santrallerinin anahtarlama devrelerini
sayısal mantıkla tasarlama problemi üzerinde çalışırken, 1953'te \emph{The Map Method for
Synthesis of Combinational Logic Circuits} (Kombinasyonel Mantık Devrelerinin Sentezi İçin
Harita Yöntemi) başlıklı makalesini yayımladı. Bu makalede tanıttığı görsel yöntem, kısa
sürede mühendislik camiasında onun adıyla anılır oldu. Karnaugh, Bell Laboratuvarları'ndaki
yıllarında darbe kodu modülasyonu (PCM) ve manyetik mantık devreleri üzerine başka
patentler de aldı; kariyerinin ikinci yarısını IBM'in Thomas J. Watson Araştırma
Merkezi'nde sürdürdü ve 1976'da IEEE Fellow seçildi. 2022'de, 98 yaşında hayatını
kaybetti.

\textbf{Kapı seviyesinde sadeleştirme}\index{kapı seviyesinde sadeleştirme} (gate-level
minimization), bir sayısal devreyi tanımlayan Bool fonksiyonlarının, mümkün olan en az kapı
ve en az girişle gerçeklenecek biçimde optimal bir kapı-seviyesi karşılığını bulma
görevidir. K-map bu göreve, doğruluk tablosunun görsel/
resimsel bir biçimi gibi yaklaşır: haritadaki her kare, fonksiyonun bir minterm'ini temsil
eder, ve fonksiyon içinde yer alan minterm'lerin kapladığı alan haritada doğrudan
\emph{görülebilir}. Bitişik karelerin oluşturduğu desenler tanınarak, aynı fonksiyonun
alternatif cebirsel ifadeleri sistematik biçimde türetilebilir; bunlar arasından terim ve
literal sayısı en az olan seçilir. Bir sonraki bölümde göreceğimiz gibi, K-map yöntemiyle
üretilen sadeleştirilmiş ifade her zaman iki standart biçimden birindedir: toplamlar-çarpımı
(sum of products, SOP) ya da çarpımlar-toplamı (product of sums, POS).

K-map'in bir sınırı vardır: kare sayısı değişken sayısıyla katlanarak arttığından ($2^n$),
yöntem elle yalnızca birkaç değişkene (pratikte 4-6) kadar rahat uygulanabilir; daha fazla
girişi olan gerçek devrelerde bu iş günümüzde bilgisayar tabanlı mantık sentezi araçlarına
bırakılır. Yine de bu araçların arkasındaki matematiği anlamak, hem küçük devreleri elle
tasarlayabilmek hem de araçların ürettiği sonucu değerlendirebilmek için vazgeçilmezdir.
Bu yüzden bölümün sonunda, K-map'in elle çizilemeyecek kadar çok değişkenli durumlar için
tablo tabanlı bir alternatifini, Quine-McCluskey algoritmasını da göreceğiz. Donanım
tanımlama dillerinin (Verilog/VHDL) bu sadeleştirilmiş ifadeleri gerçek devre modüllerine
dönüştürmede nasıl kullanıldığını ise, kombinasyonel devreleri (toplayıcı, çoklayıcı,
kod çözücü) işleyeceğimiz bir sonraki bölüme bırakıyoruz.

\section{Harita Yöntemi: İki ve Üç Değişkenli K-map}
\label{sec:kmap-yontemi}

\subsection{Harita Yöntemine Genel Bakış}
\label{sec:kmap-genel-bakis}

Bir K-map, karelerden oluşan bir diyagramdır; sadeleştirilecek fonksiyonun her bir
minterm'i için haritada tam olarak bir kare ayrılmıştır. Herhangi bir Bool fonksiyonu
minterm'lerin toplamı biçiminde yazılabildiğinden, fonksiyon haritada, kendisine ait
minterm'lerin kapladığı karelerin oluşturduğu alan olarak görsel biçimde tanınır. Bir
fonksiyonu haritaya yerleştirmenin yolu basittir: fonksiyona ait olan her minterm'in
karesine $1$ yazılır (ya da kare gölgelenir), fonksiyonda yer almayan minterm'lerin
kareleri $0$ bırakılır (veya boş bırakılır). Haritayı doğruluk tablosunun iki boyutlu,
\emph{resimsel} bir biçimi gibi düşünebilirsiniz: doğruluk tablosundaki her satır, haritada
bir kareye karşılık gelir; farkı, haritada birbirine \emph{komşu} (bitişik) kareler her
zaman birbirine algebraik olarak da yakın satırlardır. Bu bölümde önce iki ve üç değişkenli
haritaları, bir sonraki bölümde ise dört değişkenli haritayı inceleyeceğiz.

\subsection{İki Değişkenli K-map}
\label{sec:kmap-iki-degiskenli}

İki değişkenli bir fonksiyonun $4$ minterm'i vardır ($m_0,m_1,m_2,m_3$); bu yüzden harita
$2\times2=4$ kareden oluşur. Şekil~\ref{fig:kmap-iki-yapi}'de gösterildiği gibi, satırlar
$x$ değişkenine, sütunlar $y$ değişkenine ayrılmıştır: üst satır $x=0$ ($x'$), alt satır
$x=1$ (düz $x$) durumunu; sol sütun $y=0$ ($y'$), sağ sütun $y=1$ (düz $y$) durumunu
temsil eder. Her karenin hangi minterm'e karşılık geldiği, o karenin bulunduğu satır ve
sütunun kesişiminden okunur.

\begin{figure}[H]
\centering
\karnoiki{x}{y}{$m_0$,$m_1$,$m_2$,$m_3$}
\caption{İki değişkenli K-map'in yapısı ve minterm numaraları.}
\label{fig:kmap-iki-yapi}
\end{figure}

\begin{ornek}[İki Değişkenli K-map'te Fonksiyon Gösterimi]
$F_1=xy$ ve $F_2=x+y$ fonksiyonlarını haritada gösterelim. $F_1=xy$ yalnızca $m_3$'e
($x=1,y=1$) eşittir, bu yüzden haritada tek bir kareye $1$ yazılır. $F_2=x+y$ ise
$m_1+m_2+m_3$ toplamına eşittir (yani $x$ ya da $y$'nin en az biri $1$ olduğu her durum),
bu yüzden üç kareye $1$ yazılır.

\begin{center}
\karnoiki{x}{y}{0,0,0,1}\hspace{2cm}\karnoiki{x}{y}{0,1,1,1}
\end{center}
\captionof{figure}{Solda $F_1=xy$ ($m_3$), sağda $F_2=x+y$ ($m_1+m_2+m_3$).}
\label{fig:kmap-iki-ornek}
\end{ornek}

\subsection{Üç Değişkenli K-map}
\label{sec:kmap-uc-degiskenli}

Üç değişkenli bir fonksiyonun $8$ minterm'i vardır; harita da buna uygun olarak $8$
kareden oluşur (Şekil~\ref{fig:kmap-uc-yapi}). Satırlar yine $x$ değişkenine ayrılmıştır
(üst $x=0$, alt $x=1$), ama sütunlar artık \emph{iki} değişkeni ($y$ ve $z$) birlikte
temsil eder. Buradaki kritik nokta, sütunların $00,01,10,11$ gibi düz ikili sırayla
DEĞİL, $00,01,11,10$ sırasıyla dizilmiş olmasıdır: bu, Bölüm~\ref{chap:sayi-sistemleri}'de
gördüğümüz \textbf{Gray kodu}\index{Gray kodu!K-map'te kullanımı}
(bkz.~Bölüm~\ref{sec:gray-kodu}) sırasının ta kendisidir. Bu sıralama rastgele değildir:
Gray kodunun tanımlayıcı özelliği bir sonraki sütuna geçerken yalnızca \emph{tek bir}
bitin değişmesidir, ve K-map'in bütün gücü de bu özelliğe dayanır.

\begin{figure}[H]
\centering
\karnouc{x}{yz}{$m_0$,$m_1$,$m_3$,$m_2$,$m_4$,$m_5$,$m_7$,$m_6$}
\caption{Üç değişkenli K-map'in yapısı ve minterm numaraları. Sütunlar $yz$'nin Gray
kodu sırasıyla ($00,01,11,10$) dizilmiştir.}
\label{fig:kmap-uc-yapi}
\end{figure}

\begin{tanim}{Bitişik Kareler:}
Bir K-map'te iki kare, yalnızca \textbf{tek bir değişkende} farklıysa (o değişken bir
karede düz, diğerinde tümlü ise) ve geri kalan tüm değişkenler aynıysa \emph{bitişik}
(adjacent) kabul edilir. Bitişiklik yalnızca yatay ve dikey komşuluk için geçerlidir;
köşegen (diyagonal) komşu kareler bitişik SAYILMAZ.
\end{tanim}

Gray kodu sıralamasının doğrudan sonucu şudur: haritada fiziksel olarak yan yana duran
iki kare, her zaman cebirsel olarak da bitişiktir. Bu, bitişik iki minterm'in toplamının
her zaman tek bir çarpım terimine indirgenebileceği anlamına gelir; çünkü Bool
cebirinin postülatlarından, farklı olan tek değişken $y+y'=1$ olarak sadeleşir. Örneğin
$m_5=xy'z$ ve $m_7=xyz$ karelerini ele alalım (haritada aynı sütunda, farklı satırlarda
bulunurlar, yalnızca $y$ değişkeninde farklıdırlar):
\[
m_5+m_7 = xy'z+xyz = xz(y'+y) = xz
\]
İki bitişik kare toplandığında, aralarında farklı olan tek değişken kayboldu ve $3$
literallik bir ifadeden $2$ literallik bir ifadeye indik.

\begin{bilginot}[Çevresel Bitişiklik: Haritanın Kenarları Birbirine Değer]
Bitişiklik tanımı ilk bakışta göründüğünden daha geniştir. Şekil~\ref{fig:kmap-uc-yapi}'de
$m_0$ (en soldaki sütun) ile $m_2$ (en sağdaki sütun) fiziksel olarak yan yana
\emph{çizilmemiştir}, ama minterm tanımlarına bakıldığında $m_0=x'y'z'$ ile $m_2=x'yz'$
yalnızca $y$ değişkeninde farklıdır, yani tanım gereği bitişiktirler:
\[
m_0+m_2 = x'y'z'+x'yz' = x'z'(y'+y) = x'z'
\]
Aynı durum $m_4$ ile $m_6$ için de geçerlidir ($m_4+m_6=xz'$). Bunun nedeni, sütun
sıralamasının Gray kodunu takip etmesidir: ilk sütun ($00$) ile son sütun ($10$) da
yalnızca tek bir bitte ($y$) farklıdır, tıpkı ardışık iki sütun gibi. Pratik sonucu
şudur: haritayı, sağ ve sol kenarları birbirine yapışık bir \textbf{silindir} üzerine
sarılmış gibi düşünmek gerekir. Dört değişkenli haritada (bir sonraki bölüm) bu kural üst
ve alt kenarlar için de geçerli olacak.
\end{bilginot}

\begin{bilginot}[Grup Büyüklüğü Neden Mutlaka İkinin Kuvveti Olmalı?]
Bitişik kareleri ikişer ikişer birleştirme fikri, dörderli, sekizerli gruplara da
genelleşir; ama grup büyüklüğü her zaman $2$'nin bir kuvveti ($1,2,4,8,\ldots$) olmalıdır,
çünkü her birleştirme adımı bir $y+y'=1$ türü sadeleşmesine dayanır ve bu sadeleşme
karelerin sayısını her seferinde ikiye katlar. Üç değişkenli bir haritada:
\begin{itemize}[nosep, leftmargin=*]
    \item $1$ kare $\to$ $3$ literallik bir terim verir (sadeleşme yok).
    \item $2$ bitişik kare $\to$ $2$ literallik bir terim verir (bir değişken düşer).
    \item $4$ bitişik kare $\to$ $1$ literallik bir terim verir (iki değişken düşer).
    \item $8$ kare (haritanın tamamı) $\to$ fonksiyon sabit $1$'e eşittir (üç değişken de düşer).
\end{itemize}
Genel kural: $2^k$ bitişik kareden oluşan bir grup, $(n-k)$ literallik bir terime karşılık
gelir ($n$ = toplam değişken sayısı).
\end{bilginot}

\subsection{Örnekler ve Uygulama}
\label{sec:kmap-uc-degiskenli-ornekler}

\begin{ornek}[İki Ayrı İkili Grup]
\label{orn:kmap-a}
$F(x,y,z)=\Sigma(2,3,4,5)$ fonksiyonunu sadeleştirelim. Önce fonksiyona ait dört
minterm'in karesine $1$ yazılır (Şekil~\ref{fig:kmap-ornek-a}). $m_2$ ve $m_3$ bitişiktir
(aynı satır, komşu sütun) ve $x'y$ terimine indirgenir; $m_4$ ve $m_5$ de bitişiktir
(aynı satır, komşu sütun) ve $xy'$ terimine indirgenir. İki grubun birleşimi fonksiyonun
tamamını kapsadığından:
\[
F = x'y + xy'
\]

\begin{figure}[H]
\centering
\begin{tikzpicture}[scale=1.0, font=\small]
\fill[blokmavi] (2,1) rectangle (4,2);
\fill[blokturuncu] (0,0) rectangle (2,1);
\draw (0,0) grid (4,2);
\node[above] at (0.5,2) {$00$};
\node[above] at (1.5,2) {$01$};
\node[above] at (2.5,2) {$11$};
\node[above] at (3.5,2) {$10$};
\node[left] at (0,1.5) {$0$};
\node[left] at (0,0.5) {$1$};
\node[above left] at (0,2) {\scriptsize $x\backslash yz$};
\node at (0.5,1.5) {$0$}; \node at (1.5,1.5) {$0$};
\node at (2.5,1.5) {$1$}; \node at (3.5,1.5) {$1$};
\node at (0.5,0.5) {$1$}; \node at (1.5,0.5) {$1$};
\node at (2.5,0.5) {$0$}; \node at (3.5,0.5) {$0$};
\end{tikzpicture}
\caption{$F(x,y,z)=\Sigma(2,3,4,5)=x'y+xy'$.}
\label{fig:kmap-ornek-a}
\end{figure}
\end{ornek}

\begin{ornek}[Çevresel Bitişikliğin Kullanımı]
\label{orn:kmap-b}
$F(x,y,z)=\Sigma(3,4,6,7)$ fonksiyonunu sadeleştirelim (Şekil~\ref{fig:kmap-ornek-b}).
$m_3$ ve $m_7$ aynı sütunda, farklı satırlardadır; birleşince $yz$ terimini verir. $m_4$
ve $m_6$ ise haritanın en soldaki ve en sağdaki sütunlarında bulunur. Yukarıda
gördüğümüz çevresel bitişiklik kuralı gereği bunlar da bitişiktir ve $xz'$ terimini
verir. Sonuç:
\[
F = yz + xz'
\]

\begin{figure}[H]
\centering
\begin{tikzpicture}[scale=1.0, font=\small]
\fill[blokmavi] (2,0) rectangle (3,2);
\fill[blokturuncu] (0,0) rectangle (1,1);
\fill[blokturuncu] (3,0) rectangle (4,1);
\draw (0,0) grid (4,2);
\node[above] at (0.5,2) {$00$};
\node[above] at (1.5,2) {$01$};
\node[above] at (2.5,2) {$11$};
\node[above] at (3.5,2) {$10$};
\node[left] at (0,1.5) {$0$};
\node[left] at (0,0.5) {$1$};
\node[above left] at (0,2) {\scriptsize $x\backslash yz$};
\node at (0.5,1.5) {$0$}; \node at (1.5,1.5) {$0$};
\node at (2.5,1.5) {$1$}; \node at (3.5,1.5) {$0$};
\node at (0.5,0.5) {$1$}; \node at (1.5,0.5) {$0$};
\node at (2.5,0.5) {$1$}; \node at (3.5,0.5) {$1$};
\end{tikzpicture}
\caption{$F(x,y,z)=\Sigma(3,4,6,7)=yz+xz'$. Turuncu renkli iki kare ($m_4$ ve $m_6$)
haritanın zıt kenarlarında olsa da çevresel bitişiklik sayesinde tek bir grup oluşturur.}
\label{fig:kmap-ornek-b}
\end{figure}
\end{ornek}

\begin{ornek}[Dörtlü Grup ve Kare Tekrar Kullanımı]
\label{orn:kmap-c}
$F(x,y,z)=\Sigma(0,2,4,5,6)$ fonksiyonunu sadeleştirelim (Şekil~\ref{fig:kmap-ornek-c}).
$m_0,m_2,m_4,m_6$ kareleri hep birlikte $4$'lü bir grup oluşturur (haritanın en soldaki
ve en sağdaki sütunlarının tamamı, çevresel bitişiklikle birleşir): bu dört karede ortak
olan tek şey $z=0$ olmasıdır, bu yüzden grup $z'$ terimine indirgenir. Geriye tek başına
kalan $m_5$ ise, daha önce $z'$ grubunda kullanılmış olan $m_4$ ile tekrar birleştirilebilir
(bir kare birden fazla grupta kullanılabilir): $m_4+m_5=xy'$. Sonuç:
\[
F = z' + xy'
\]

\begin{figure}[H]
\centering
\begin{tikzpicture}[scale=1.0, font=\small]
\fill[blokmavi] (0,0) rectangle (1,2);
\fill[blokmavi] (3,0) rectangle (4,2);
\draw[very thick, vurguturuncu, rounded corners=2pt] (0.12,0.12) rectangle (2-0.12,1-0.12);
\draw (0,0) grid (4,2);
\node[above] at (0.5,2) {$00$};
\node[above] at (1.5,2) {$01$};
\node[above] at (2.5,2) {$11$};
\node[above] at (3.5,2) {$10$};
\node[left] at (0,1.5) {$0$};
\node[left] at (0,0.5) {$1$};
\node[above left] at (0,2) {\scriptsize $x\backslash yz$};
\node at (0.5,1.5) {$1$}; \node at (1.5,1.5) {$0$};
\node at (2.5,1.5) {$0$}; \node at (3.5,1.5) {$1$};
\node at (0.5,0.5) {$1$}; \node at (1.5,0.5) {$1$};
\node at (2.5,0.5) {$0$}; \node at (3.5,0.5) {$1$};
\end{tikzpicture}
\caption{$F(x,y,z)=\Sigma(0,2,4,5,6)=z'+xy'$. Mavi dörtlü grup $z'$ terimini, turuncu
çerçeveli ikili grup (biri mavi grupla ortak) $xy'$ terimini verir.}
\label{fig:kmap-ornek-c}
\end{figure}
\end{ornek}

\begin{ornek}[Kanonik Olmayan Bir İfadeden Minterm Okumak]
\label{orn:kmap-d}
Şimdiye kadarki örneklerde fonksiyon zaten minterm toplamı ($\Sigma$) biçiminde
verilmişti. Ama bir fonksiyon toplamlar-çarpımı (SOP) biçiminde olsa bile, doğrudan
minterm toplamı olmayabilir, yani terimlerinden bazıları birden fazla değişken
içerebilir. Böyle bir ifadeyi de haritaya yerleştirebiliriz, her çarpım terimini
kapsadığı kareler kümesine genişleterek. Örneğin:
\[
F(A,B,C) = AC + A'B + ABC' + B'C
\]
$AC$ terimi ($A=1,C=1$, $B$ serbest) $m_5$ ve $m_7$'yi kapsar. $A'B$ terimi ($A=0,B=1$,
$C$ serbest) $m_2$ ve $m_3$'ü kapsar. $ABC'$ terimi tek başına $m_6$'yı kapsar. $B'C$
terimi ($B=0,C=1$, $A$ serbest) $m_1$ ve $m_5$'i kapsar ($m_5$ zaten işaretliydi, tekrar
işaretlemeye gerek yok). Toplamda işaretlenen kareler $\{1,2,3,5,6,7\}$'dir, yani
$F(A,B,C)=\Sigma(1,2,3,5,6,7)$. Şimdi bu haritayı yeniden gruplarsak (Şekil
\ref{fig:kmap-ornek-d}), çok daha sade bir sonuç ortaya çıkar: sağdaki iki sütun
($B=1$ olan tüm kareler) bir grup oluşturup $B$ terimini, ortadaki iki sütun ($C=1$
olan tüm kareler) başka bir grup oluşturup $C$ terimini verir:
\[
F = B + C
\]
Bu, orijinal dört terimli ifadenin yalnızca iki tek-literallik terime indiği çarpıcı bir
örnektir. Cebirsel yoldan bu sonuca ulaşmak çok daha fazla adım gerektirirdi.

\begin{figure}[H]
\centering
\begin{tikzpicture}[scale=1.0, font=\small]
\fill[blokmavi] (2,0) rectangle (4,2);
\draw[very thick, vurguturuncu, rounded corners=2pt] (1.12,0.12) rectangle (3-0.12,2-0.12);
\draw (0,0) grid (4,2);
\node[above] at (0.5,2) {$00$};
\node[above] at (1.5,2) {$01$};
\node[above] at (2.5,2) {$11$};
\node[above] at (3.5,2) {$10$};
\node[left] at (0,1.5) {$0$};
\node[left] at (0,0.5) {$1$};
\node[above left] at (0,2) {\scriptsize $A\backslash BC$};
\node at (0.5,1.5) {$0$}; \node at (1.5,1.5) {$1$};
\node at (2.5,1.5) {$1$}; \node at (3.5,1.5) {$1$};
\node at (0.5,0.5) {$0$}; \node at (1.5,0.5) {$1$};
\node at (2.5,0.5) {$1$}; \node at (3.5,0.5) {$1$};
\end{tikzpicture}
\caption{$F(A,B,C)=\Sigma(1,2,3,5,6,7)=B+C$. Mavi grup ($B=1$) $B$ terimini, turuncu
çerçeveli grup ($C=1$, ortadaki iki kareyi mavi grupla paylaşır) $C$ terimini verir.}
\label{fig:kmap-ornek-d}
\end{figure}
\end{ornek}

\begin{alistirma}[\kolay]
$F(x,y,z)=\Sigma(0,1,6,7)$ fonksiyonunu sadeleştiriniz.
\end{alistirma}

\noindent\textbf{Çözüm:} $m_0$ ve $m_1$ bitişiktir (aynı satır, komşu sütun), $z$
değişkeni düşer ve $x'y'$ terimi kalır. $m_6$ ve $m_7$ de bitişiktir, yine $z$ düşer ve
$xy$ terimi kalır. Sonuç: $F = x'y' + xy$.

\begin{alistirma}[\kolay]
$F(x,y,z)=\Sigma(0,1,2,5)$ fonksiyonunu sadeleştiriniz.
\end{alistirma}

\noindent\textbf{Çözüm:} $m_0$ ve $m_2$ çevresel bitişiklikle birleşir ($y$ düşer),
$x'z'$ terimini verir. $m_1$ ve $m_5$ bitişiktir ($x$ düşer), $y'z$ terimini verir. İki
grup birlikte $\{0,2,1,5\}$'in tamamını kapsar. Sonuç: $F = x'z' + y'z$.

\begin{alistirma}[\orta]
$F(x,y,z)=\Sigma(0,2,3,4,6)$ fonksiyonunu sadeleştiriniz.
\end{alistirma}

\noindent\textbf{Çözüm:} $m_0,m_2,m_4,m_6$ çevresel bitişiklikle dörtlü bir grup oluşturur
($z=0$ ortak), $z'$ terimini verir. Geriye kalan $m_3$, daha önce kullanılan $m_2$ ile
tekrar birleştirilebilir ($x$ düşer), $x'y$ terimini verir. Sonuç: $F = z' + x'y$.

\begin{alistirma}[\zor]
$F(A,B,C)=A'BC+AB'+AC'+B'C$ ifadesi için (a) fonksiyonu minterm toplamı olarak yazınız,
(b) minimal toplamlar-çarpımı (SOP) ifadesini bulunuz.
\end{alistirma}

\noindent\textbf{Çözüm:} (a) $A'BC$ terimi $m_3$'ü, $AB'$ terimi $m_4$ ve $m_5$'i, $AC'$
terimi $m_4$ ve $m_6$'yı, $B'C$ terimi $m_1$ ve $m_5$'i kapsar. Birleşim
$\{1,3,4,5,6\}$'dır, yani $F(A,B,C)=\Sigma(1,3,4,5,6)$.
(b) Haritaya yerleştirilip yeniden gruplandığında: $m_4$ ve $m_5$ bitişiktir ve $AB'$
terimini verir; $m_4$ ve $m_6$ bitişiktir ve $AC'$ terimini verir; $m_1$ ve $m_3$
bitişiktir ve $A'C$ terimini verir. Üç grup birlikte $\{1,3,4,5,6\}$'nın tamamını
kapsar ($m_4$ iki grupta paylaşılır). Sonuç: $F = AB' + AC' + A'C$.
